\section{Сравнение существующих подходов и анализ аналогов}

\subsection{Сравнение существующих подходов}

\subsubsection{Современные направления в Text-to-SQL}

Современные подходы к решению задачи \textbf{Text-to-SQL} представляют собой переход от классических нейросетевых архитектур к сложным многоагентным системам и методам глубокого поиска, использующим возможности \textbf{больших языковых моделей (LLM)}.  \cite{hong2025survey} Текущее состояние области характеризуется разделением на две доминирующие парадигмы: \textbf{In-Context Learning (ICL)}, опирающуюся на продвинутый промпт-инжиниринг проприетарных моделей, и \textbf{Fine-tuning (FT)} открытых моделей для специфических диалектов и структур данных. \cite{li2024bird}

Одним из наиболее перспективных направлений является интеграция методов \textbf{поиска по дереву Монте-Карло (MCTS)}, реализованная в подходе \textbf{Alpha-SQL}. \cite{li2025alphasql} Этот метод формулирует генерацию SQL как задачу \textbf{итеративного поиска} в пространстве действий, где LLM выступает в роли модели действий, генерируя пошаговые рассуждения (Chain-of-Thought) для каждого узла дерева. Alpha-SQL использует самообучаемую функцию вознаграждения, основанную на \textbf{согласованности результатов выполнения} (self-consistency), что позволяет моделям среднего размера (например, 32B) превосходить такие гиганты, как GPT-4o, на сложных бенчмарках.

Для решения проблем \textbf{сложного математического анализа} и навигации по запутанным схемам предложен подход \textbf{SteinerSQL}. \cite{mao2025steinersql} В рамках этого фреймворка задача поиска оптимального пути соединения таблиц (join path) сводится к классической задаче \textbf{дерева Штейнера на графе схемы}. SteinerSQL объединяет математическую декомпозицию вопроса с детерминированными алгоритмами на графах, что гарантирует сохранение вычислительных зависимостей и минимизирует «шум» от избыточных соединений.

Проблема \textbf{надежности и галлюцинаций} в многоагентных системах эффективно решается с помощью технологии \textbf{consistency alignment} (выравнивание согласованности), представленной в \textbf{OpenSearch-SQL}. \cite{xie2025opensearch} Этот подход внедряет специализированный агент выравнивания, который гарантирует, что входные и выходные данные различных модулей (препроцессинг, экстракция, генерация) логически соответствуют друг другу, предотвращая накопление ошибок в цепочке агентов. Дополнительно в систему внедрен промежуточный язык \textbf{SQL-Like}, который упрощает структуру запроса перед генерацией финального синтаксиса.

В области \textbf{обучения с подкреплением} выделяется семейство моделей \textbf{Think2SQL}, использующее парадигму \textbf{RLVR} (Reinforcement Learning with Verifiable Rewards). \cite{papicchio2026think2sql} Вместо субъективных человеческих предпочтений модели обучаются на основе \textbf{проверяемых вознаграждений} от движка базы данных, используя плотные функции награды (Cell Precision, Cell Recall) для обеспечения стабильной конвергенции даже небольших моделей. Это позволяет модели «открывать» новые пути рассуждения, которые отсутствовали в исходных наборах данных для дообучения.

Для корпоративных систем, требующих \textbf{статистических гарантий корректности}, разработан фреймворк \textbf{RTS++}. \cite{chen2026rts} Он вводит понятие \textbf{execution entropy} (энтропия выполнения) для количественной оценки неопределенности модели и использует \textbf{конформное прогнозирование} (conformal prediction) для принятия решения о выполнении запроса или его передаче человеку на проверку. Такой подход критически важен в средах, где «тихие» ошибки (синтаксически верные, но семантически неверные запросы) недопустимы.

Наконец, стремление к \textbf{вычислительной эффективности} и защите данных привело к созданию \textbf{LitE-SQL} — легковесного фреймворка, демонстрирующего, что модели с 7B параметров могут достигать точности, сопоставимой с моделями в 175B–200B параметров. \cite{piao2025litesql} Успех этого подхода обеспечивается \textbf{векторным связыванием схемы} (Schema Linking) на основе контрастивного обучения и двухэтапным дообучением, что делает систему пригодной для локального развертывания.

\subsubsection{Современные направления в RAG для баз данных}

Современные направления в области \textbf{Retrieval-Augmented Generation (RAG)} для баз данных и их интеграция в задачи \textbf{Text-to-SQL} представляют собой переход от простых поисковых систем к сложным «Knowledge Runtime» средам. \cite{piao2025litesql} В 2026 году RAG уже не рассматривается как вспомогательный компонент, а становится ключевым слоем инженерии поиска, обеспечивающим актуальность знаний, соблюдение прав доступа и минимизацию стоимости эксплуатации за счет использования \textbf{малых языковых моделей (SLM)}. \cite{kmoseenk2026return}

Одним из основных направлений интеграции RAG в Text-to-SQL является \textbf{связывание схемы (Schema Linking)} на основе векторных баз данных. \cite{piao2025litesql} В отличие от ранних методов, современные системы, такие как \textbf{LitE-SQL}, хранят предобученные эмбеддинги метаданных схемы (названия таблиц, описания столбцов) в векторном хранилище и извлекают только релевантные элементы в ответ на запрос пользователя. \cite{piao2025litesql, piao2025arxiv} Это позволяет эффективно работать с огромными корпоративными схемами, содержащими тысячи столбцов, не переполняя контекстное окно модели и минимизируя шум от нерелевантных данных. \cite{ji2026benchmarks, piao2025litesql}

\textbf{Гибридный поиск (Hybrid Search)} и \textbf{реранкинг (Reranking)} признаны наиболее эффективными способами повышения точности извлечения. \cite{akzhankalimatov2026rag, kmoseenk2026return} Подход сочетает семантический векторный поиск, улавливающий общие смыслы, с лексическим поиском (BM25) для точного сопоставления технических терминов, идентификаторов и кодов ошибок. \cite{lin2025higress} Использование алгоритма \textbf{Reciprocal Rank Fusion (RRF)} для объединения результатов и последующее использование \textbf{кросс-энкодеров} для переранжирования топ-результатов позволяет повысить качество извлечения на 15–40\%. \cite{lin2025higress}

Развитие \textbf{агентного RAG (Agentic RAG)} привело к созданию модульных систем, где специализированные агенты (например, \textbf{Schema-RAG-Agent}) автономно планируют свои действия. \cite{chudinov2025agentic, xie2025opensearch} Агенты могут переформулировать запрос пользователя, решать, к какому источнику обратиться (документация, SQL, Slack), и проверять релевантность найденного контекста перед генерацией. \cite{akzhankalimatov2026rag, xie2025opensearch} Фреймворки типа \textbf{Corrective RAG (CRAG)} и \textbf{Self-RAG} добавляют уровень самокоррекции: если извлеченный контекст признан нерелевантным или амбивалентным, система инициирует повторный поиск или обращается к внешним инструментам (web search). \cite{akzhankalimatov2026rag, lin2025higress}

Инновационной парадигмой является \textbf{Table-Augmented Generation (TAG)}, которая объединяет Text-to-SQL и RAG для обработки запросов, требующих одновременно семантического анализа (мировых знаний) и точных вычислений в БД. \cite{biswal2024tag} Параллельно предложена концепция \textbf{Dual-Retrieval-Augmented-Generation (DRAG)}, формализующая рабочий процесс двойного извлечения схем таблиц и документов знаний перед генерацией SQL-запроса. \cite{ji2026benchmarks}

Для обеспечения надежности в корпоративных системах внедряются механизмы \textbf{выравнивания согласованности (Consistency Alignment)}, такие как в \textbf{OpenSearch-SQL}, и использование «энтропии выполнения» для оценки уверенности модели. \cite{xie2025opensearch, chen2026rts} Это превращает Text-to-SQL из вероятностной генерации в надежный сервис, способный воздержаться от ответа при высокой неопределенности. \cite{chen2026rts}


\subsection{Сравнение конкурентов}

В рамках разработки интеллектуального модуля был проведён сравнительный анализ программных решений, которые позволяют автоматизировать взаимодействие с реляционными базами данных на естественном языке. В качестве аналогов рассматривались наиболее востребованные инструменты в области \texttt{Text-to-SQL} и бизнес-аналитики — Vanna.AI, LangChain и Metabase.

\textbf{Vanna.AI} представляет собой специализированный Python-фреймворк с открытым исходным кодом, построенный на технологии \texttt{RAG}. Система обучается на метаданных, \texttt{DDL}-запросах и примерах пользовательских запросов, что обеспечивает высокую точность генерации для специфических схем баз данных. Однако, как показал анализ, отсутствует надёжный механизм самокоррекции: если сгенерированный \texttt{SQL}-запрос вызывает ошибку в \texttt{СУБД}, пользователю приходится править его вручную. К тому же возможности визуализации ограничены стандартными библиотеками и требуют дополнительного кода.

\textbf{LangChain} с его модулем \texttt{SQL Database Chain} — один из самых популярных фреймворков для создания приложений на базе больших языковых моделей. \cite{official_langchain} Он позволяет создавать цепочки и агентов, которые могут обращаться к схеме базы данных в реальном времени, и радует гибкостью настройки и обилием интеграций. Но у этого решения есть слабые места: стандартные цепочки LangChain для \texttt{SQL} начинают пробуксовывать на очень больших схемах (например, на демонстрационной базе «Авиаперевозки» с её множеством связей). Модель часто пытается передать в контекст всю схему целиком, что приводит к ошибкам или перерасходу токенов. Механизмы автоматического исправления ошибочных запросов в стандартных цепочках развиты слабо.

\textbf{Metabase} — это зрелая платформа для бизнес-аналитики, которая недавно обзавелась \texttt{AI}-помощником для генерации вопросов к данным. Она ориентирована на конечных бизнес-пользователей и создание дашбордов, предлагая отличную визуализацию «из коробки» и удобный интерфейс. Но Metabase — преимущественно облачное или серверное решение, что создаёт риски для конфиденциальности данных. Его \texttt{AI}-функционал слабо адаптирован для глубокой работы с \texttt{RAG}-системами на стороне клиента, а механизм самокоррекции практически отсутствует.

Спроектированный в рамках данной дипломной работы интеллектуальный модуль призван устранить перечисленные недостатки. Ключевое отличие — объединение трёх важных факторов в едином контуре. Во-первых, это автономность и приватность: использование локальных языковых моделей гарантирует, что данные авиапредприятия никогда не покинут внутренний периметр. Во-вторых, механизм самокоррекции: модуль реализует агентский подход, при котором в случае ошибки \texttt{СУБД} агент анализирует текст ошибки и автоматически пересобирает \texttt{SQL}-запрос до получения валидного результата. \cite{chaturvedi2025sql} В-третьих, глубокая проработка \texttt{RAG} для сложных схем. \cite{arxiv2410rag} В отличие от конкурентов, система оптимизирована для работы со специфическими связями базы данных Postgres Pro «Авиаперевозки» и использует векторное хранилище, чтобы извлекать только действительно нужные фрагменты схемы.

Сводный анализ функциональных возможностей исследуемых решений представлен в таблице~\ref{table:comparison}.

\begin{table}[h!]
    \centering
    \small   % уменьшаем шрифт (можно \footnotesize)
    \caption{Сравнение разрабатываемого решения с существующими инструментами}
    \label{table:comparison}
    % Устанавливаем относительную ширину столбцов:
    % первый столбец (функционал) шире, остальные — равномерно
    \begin{tabularx}{\textwidth}{|>{\hsize=1.2\hsize\raggedright\arraybackslash}X|
                                  |>{\hsize=0.95\hsize\centering\arraybackslash}X|
                                  |>{\hsize=0.95\hsize\centering\arraybackslash}X|
                                  |>{\hsize=0.95\hsize\centering\arraybackslash}X|
                                  |>{\hsize=0.95\hsize\centering\arraybackslash}X|}
    \hline
    \textbf{Функционал} & \textbf{Vanna.AI} & \textbf{LangChain SQL} & \textbf{Metabase} & \textbf{Мой продукт} \\
    \hline
    Генерация SQL (Text-to-SQL) & + & + & + & + \\
    \hline
    RAG для схемы БД & + & Ограничено & -- & + \\
    \hline
    Локальный запуск (Privacy) & + & + & -- & + \\
    \hline
    Механизм самокоррекции (Self-Heal) & -- & -- & -- & + \\
    \hline
    Авто-визуализация данных & $\pm$ & -- & + & Тесты \\
    \hline
    \end{tabularx}
\end{table} 

